\documentclass[12pt, letterpaper]{article}

\usepackage[utf8]{inputenc}
\usepackage[T1]{fontenc}
\usepackage[spanish]{babel}
\usepackage{geometry}
\usepackage{amsmath, amssymb, amsthm}
\usepackage{booktabs}
\usepackage{array}
\usepackage{microtype}
\usepackage{setspace}
\usepackage{parskip}
\usepackage{titlesec}
\usepackage{fancyhdr}
\usepackage{enumitem}
\usepackage{bm}
\usepackage{mathtools}
\usepackage{tcolorbox}
\usepackage{hyperref}

\tcbuselibrary{skins, breakable}

\geometry{top=2.8cm, bottom=3.0cm, left=3.2cm, right=3.2cm}
\setstretch{1.20}

\pagestyle{fancy}
\fancyhf{}
\fancyhead[C]{\small\textsc{Econometria · Gujarati --- Ejercicio 10.10}}
\fancyfoot[C]{\small\thepage}
\renewcommand{\headrulewidth}{0.4pt}
\renewcommand{\footrulewidth}{0pt}

\titleformat{\section}
  {\large\bfseries\scshape}{\thesection.}{0.6em}{}
  [\vspace{-4pt}\rule{\linewidth}{0.4pt}]
\titleformat{\subsection}{\normalsize\bfseries}{\thesubsection.}{0.5em}{}
\titleformat{\subsubsection}{\normalsize\itshape}{}{0em}{}

\newtheorem{prop}{Proposicion}
\newtheorem{lema}{Lema}
\theoremstyle{remark}
\newtheorem{obs}{Observacion}

\newtcolorbox{clave}{
  colback=white, colframe=black,
  boxrule=0.5pt, arc=3pt,
  left=8pt, right=8pt, top=6pt, bottom=6pt,
  breakable
}

\newcommand{\E}{\mathbb{E}}
\newcommand{\Var}{\operatorname{Var}}
\newcommand{\Cov}{\operatorname{Cov}}
\newcommand{\plim}{\operatorname{plim}}

\begin{document}

\begin{center}
  {\LARGE\bfseries Regresion Polinomial, Correlaciones de Orden Cero}\\[6pt]
  {\LARGE\bfseries y Multicolinealidad en la Funcion de Costo Cubica}\\[12pt]
  {\large\itshape Una distincion entre dependencia funcional y multicolinealidad lineal}\\[14pt]
  \rule{0.55\linewidth}{0.6pt}\\[8pt]
  {\normalsize Ejercicio~10.10 --- \textit{Econometria}, Gujarati \& Porter, 5.a ed.}\\[4pt]
  {\small\textsc{Derivaciones formales, geometria polinomial y sesgo de especificacion}}
\end{center}

\vspace{10pt}

\noindent\textbf{Contexto.}\quad
El ejemplo~7.4 estima una funcion de costo cubica para una empresa, donde la
variable dependiente es el costo total $Y_i$ y los regresores son el nivel de
produccion $X_i$, su cuadrado $X_i^2$ y su cubo $X_i^3$. El modelo es:

\begin{equation}
  Y_i = \beta_1 + \beta_2 X_i + \beta_3 X_i^2 + \beta_4 X_i^3 + u_i
  \label{eq:cubica}
\end{equation}

\noindent La matriz de correlaciones de orden cero reportada es:

\begin{table}[ht]
\centering
\caption{Matriz de correlaciones de orden cero entre los regresores}
\label{tab:corr}
\small
\begin{tabular}{lcccc}
\toprule
         & $X_i$  & $X_i^2$ & $X_i^3$ \\
\midrule
$X_i$    & $1.00$ & $0.9742$ & $0.9284$ \\
$X_i^2$  &        & $1.00$   & $0.9872$ \\
$X_i^3$  &        &          & $1.00$   \\
\bottomrule
\end{tabular}
\end{table}

\noindent Todas las correlaciones superan $0.92$, lo que a primera vista sugiere
multicolinealidad grave. A continuacion se analiza si esta interpretacion es correcta
y cuales son las implicaciones practicas.

\section{Inciso (a): las correlaciones altas implican multicolinealidad grave?}

\subsection{La distincion fundamental: dependencia funcional vs. colinealidad lineal}

El supuesto de no multicolinealidad del modelo de regresion clasico exige que
no exista una \textbf{relacion lineal exacta} entre los regresores en la muestra.
Formalmente, la condicion que se debe evitar es que existan escalares
$\lambda_1, \lambda_2, \lambda_3$ no todos nulos tales que:

\begin{equation}
  \lambda_2 X_i + \lambda_3 X_i^2 + \lambda_4 X_i^3 = 0
  \qquad \text{para todo } i = 1, \ldots, n
  \label{eq:colineal_exacta}
\end{equation}

En la regresion polinomial, las variables $X_i$, $X_i^2$ y $X_i^3$ son
\textbf{transformaciones no lineales de la misma variable base} $X_i$.
La relacion entre ellas es de tipo potencial, no lineal:
\begin{equation}
  X_i^2 = (X_i)^2 \neq aX_i + b \qquad \text{en general}
  \label{eq:no_lineal}
\end{equation}

\noindent Por lo tanto, la condicion~(\ref{eq:colineal_exacta}) no se cumple
para ningun conjunto de constantes no triviales, y la multicolinealidad
\emph{perfecta} no existe entre estos regresores.

\subsection{Por que las correlaciones de orden cero son altas}

La alta correlacion muestral entre $X_i$, $X_i^2$ y $X_i^3$ tiene una
explicacion aritmetica, no economica. Considere que si $X_i > 0$ para todo $i$
(como es el caso para niveles de produccion), entonces:

\begin{equation}
  r(X_i,\, X_i^2)
    = \frac{\sum x_i x_i^2}{\sqrt{\sum x_i^2 \sum (x_i^2)^2}}
    = \frac{\sum x_i^3}{\sqrt{\sum x_i^2 \sum x_i^4}}
  \label{eq:corr_x_x2}
\end{equation}

\noindent donde $x_i = X_i - \bar X$ y $x_i^2 = X_i^2 - \overline{X^2}$.
Si la distribucion de $X_i$ esta concentrada en valores positivos y no tiene
simetria alrededor de cero, los momentos de orden superior son todos positivos
y la correlacion es naturalmente alta.

Esto es un fenomeno muestral vinculado a la escala y rango de los datos, no
a una dependencia lineal entre los regresores. Cuando la produccion oscila,
por ejemplo, entre 1 y 10 unidades, los valores de $X^2$ (1 a 100) y $X^3$
(1 a 1000) son proporcionales a $X$ en esa muestra especifica, pero
globalmente la relacion es curvilinea.

\subsection{Consecuencias practicas: VIF elevado pero estimadores validos}

Aunque la multicolinealidad no es perfecta, los VIF de los regresores
polinomiales son tipicamente altos. Para el modelo cubico con las
correlaciones de la Tabla~\ref{tab:corr}:

\begin{equation}
  \text{VIF}(X_i^2) = \frac{1}{1 - R_{X_i^2 \cdot X_i, X_i^3}^2}
  \label{eq:VIF_X2}
\end{equation}

\noindent donde $R_{X_i^2 \cdot X_i, X_i^3}^2$ es el $R^2$ de la regresion
auxiliar de $X_i^2$ sobre $X_i$ y $X_i^3$. Con $r_{12} = 0.9742$ y
$r_{23} = 0.9872$, este $R^2$ es muy cercano a~1 y el VIF es grande.

Sin embargo, la alta colinealidad imperfecta solo infla los errores estandar;
\emph{no sesga los estimadores}. Si los coeficientes de~(\ref{eq:cubica}) son
estadisticamente significativos a pesar del VIF elevado, la muestra contiene
suficiente variacion independiente para identificarlos.

\begin{clave}
  \textbf{Respuesta al inciso~(a).} La afirmacion es \textbf{parcialmente
  incorrecta}. Que las correlaciones de orden cero sean altas no implica
  automaticamente multicolinealidad grave en el sentido de invalidar la
  estimacion. En regresiones polinomiales, las altas correlaciones entre
  $X$, $X^2$ y $X^3$ son un artefacto aritmetico de la escala de los
  datos, no una dependencia lineal estructural. La multicolinealidad
  perfecta no existe entre potencias de la misma variable (a menos que el
  rango de $X$ sea tan estrecho que $X^2 \approx cX + d$). La colinealidad
  imperfecta infla varianzas pero no sesga los estimadores.
\end{clave}

\section{Inciso (b): ¿eliminar \texorpdfstring{$X_i^2$}{Xi2} y \texorpdfstring{$X_i^3$}{Xi3}?}

\subsection{Respuesta}

\textbf{No}. Eliminar $X_i^2$ y $X_i^3$ basandose en la alta correlacion
seria un error metodologico grave por las siguientes razones.

\subsection{Razon 1: fundamento teorico de la funcion de costo cubica}

La teoria microeconomica del costo a corto plazo establece que:

\begin{itemize}[leftmargin=1.6em, itemsep=4pt]
  \item El costo marginal ($CM$) tiene forma de U: primero decrece
    (rendimientos crecientes de los factores variables) y luego crece
    (rendimientos decrecientes). Esto implica que el $CM$ es una funcion
    cuadratica de $X$.
  \item Integrando: el costo total debe ser una funcion cubica de $X$.
  \item El costo variable promedio ($CVP$) tambien tiene forma de U, lo que
    requiere un termino cuadratico.
\end{itemize}

Formalmente, si $CM = \partial Y/\partial X = \beta_2 + 2\beta_3 X + 3\beta_4 X^2$,
para que esta funcion tenga forma de U se necesita $\beta_3 < 0$ y $\beta_4 > 0$,
lo que requiere los tres terminos en el modelo.

\subsection{Razon 2: el criterio de omision correcto es la relevancia, no la correlacion}

La colinealidad entre regresores no es criterio valido para omitir variables.
La decision de omitir una variable debe basarse en:

\begin{enumerate}[leftmargin=1.8em, itemsep=4pt]
  \item Evidencia empirica de que su coeficiente es cero (test $t$ no significativo).
  \item Justificacion teorica de que la variable es irrelevante.
\end{enumerate}

Si los coeficientes de $X_i^2$ y $X_i^3$ son significativos (como sugiere
la teoria de costos), omitirlos introduce sesgo de especificacion independientemente
de la correlacion que tengan con $X_i$.

\subsection{Razon 3: la colinealidad se resuelve de otras formas}

Si la colinealidad entre los terminos polinomiales genera errores estandar
inaceptables, el remedio correcto es la \textbf{ortogonalizacion polinomial},
no la eliminacion de terminos. Los \textbf{polinomios ortogonales de Legendre}
transforman $X_i$, $X_i^2$, $X_i^3$ en combinaciones lineales $P_0(X)$,
$P_1(X)$, $P_2(X)$, $P_3(X)$ tales que:

\begin{equation}
  \sum_{i=1}^n P_j(X_i)\, P_k(X_i) = 0 \quad \text{para } j \neq k
  \label{eq:ortogonal}
\end{equation}

El modelo estimado con polinomios ortogonales es algebraicamente equivalente
al modelo cubico original pero sin colinealidad entre regresores, por lo que
los errores estandar son minimos y cada coeficiente se puede probar
independientemente.

\begin{clave}
  \textbf{Respuesta al inciso~(b).} No se deben eliminar $X_i^2$ y $X_i^3$.
  Son terminos teoricamente necesarios para capturar la forma en U del costo
  marginal. Si la colinealidad es un problema practico, el remedio es
  la ortogonalizacion de los terminos polinomiales, no su eliminacion.
\end{clave}

\section{Inciso (c): que sucede con el coeficiente de \texorpdfstring{$X_i$}{Xi} si se eliminan \texorpdfstring{$X_i^2$}{Xi2} y \texorpdfstring{$X_i^3$}{Xi3}}

\subsection{Derivacion formal del sesgo}

Suponga que el verdadero modelo es~(\ref{eq:cubica}) pero se estima el modelo
mal especificado:

\begin{equation}
  Y_i = \alpha_1 + \alpha_2 X_i + e_i
  \label{eq:lineal}
\end{equation}

Sea $b_{2,X^2}$ el coeficiente de la regresion de $X_i^2$ sobre $X_i$, y
$b_{2,X^3}$ el coeficiente de la regresion de $X_i^3$ sobre $X_i$ (ambas en
desviaciones). El estimador MCO $\hat\alpha_2$ del modelo~(\ref{eq:lineal}) es:

\begin{equation}
  \hat\alpha_2 = \frac{\sum x_i y_i}{\sum x_i^2}
  \label{eq:alpha2}
\end{equation}

Sustituyendo el verdadero modelo en~(\ref{eq:alpha2}):

\begin{align}
  \hat\alpha_2
  &= \frac{\sum x_i (\beta_2 x_i + \beta_3 x_i^{(2)} + \beta_4 x_i^{(3)} + u_i)}
          {\sum x_i^2}
  \notag\\[4pt]
  &= \beta_2 + \beta_3 \cdot \frac{\sum x_i x_i^{(2)}}{\sum x_i^2}
            + \beta_4 \cdot \frac{\sum x_i x_i^{(3)}}{\sum x_i^2}
            + \frac{\sum x_i u_i}{\sum x_i^2}
  \notag\\[4pt]
  &= \beta_2 + \beta_3 b_{2,X^2} + \beta_4 b_{2,X^3} + \frac{\sum x_i u_i}{\sum x_i^2}
  \label{eq:alpha2_descomp}
\end{align}

donde $x_i^{(k)} = X_i^k - \overline{X^k}$ son las desviaciones de los
terminos polinomiales. Tomando esperanza:

\begin{equation}
  \boxed{
    \E[\hat\alpha_2]
      = \beta_2
        + \underbrace{\beta_3 \cdot b_{2,X^2}}_{\text{sesgo por } X^2}
        + \underbrace{\beta_4 \cdot b_{2,X^3}}_{\text{sesgo por } X^3}
  }
  \label{eq:sesgo_total}
\end{equation}

\subsection{Signo y naturaleza del sesgo}

Para una funcion de costo cubica con la forma teorica correcta:

\begin{align}
  \beta_2 &> 0 \quad \text{(costo marginal inicial positivo)}\\
  \beta_3 &< 0 \quad \text{(tramo de rendimientos crecientes)}\\
  \beta_4 &> 0 \quad \text{(tramo de rendimientos decrecientes)}
\end{align}

Los coeficientes de las regresiones auxiliares son:

\begin{equation}
  b_{2,X^2} = r_{X,X^2} \cdot \frac{s_{X^2}}{s_X} > 0
  \label{eq:b_X2}
\end{equation}

\begin{equation}
  b_{2,X^3} = r_{X,X^3} \cdot \frac{s_{X^3}}{s_X} > 0
  \label{eq:b_X3}
\end{equation}

Los dos terminos de sesgo tienen signos opuestos:

\begin{align}
  \beta_3 b_{2,X^2} &< 0 \quad (\beta_3 < 0, \; b_{2,X^2} > 0)\\
  \beta_4 b_{2,X^3} &> 0 \quad (\beta_4 > 0, \; b_{2,X^3} > 0)
\end{align}

El sesgo neto depende de la magnitud relativa de estos dos efectos.
En general:

\begin{equation}
  \E[\hat\alpha_2] \neq \beta_2
  \label{eq:sesgo_neto}
\end{equation}

\subsection{Perdida de la interpretacion economica}

El coeficiente $\beta_2$ en el modelo completo tiene una interpretacion
precisa: es la pendiente del costo total respecto a la produccion en el
punto donde $X = 0$, es decir, el costo marginal inicial. El estimador
$\hat\alpha_2$ del modelo reducido ya no tiene esta interpretacion: es
una mezcla de los efectos lineal, cuadratico y cubico de la produccion
sobre el costo.

En particular, el modelo lineal~(\ref{eq:lineal}) implica que el costo
marginal es constante e igual a $\hat\alpha_2$, lo que contradice la
teoria del costo a corto plazo.

\subsection{Consecuencias sobre la varianza del error}

El error del modelo reducido es:

\begin{equation}
  e_i = u_i + \beta_3 X_i^2 + \beta_4 X_i^3
  \label{eq:error_reducido}
\end{equation}

(en terminos de las desviaciones de los regresores omitidos). La varianza del
error del modelo reducido es:

\begin{equation}
  \sigma_e^2
    = \sigma_u^2
    + \beta_3^2 \sigma_{X^2}^2
    + \beta_4^2 \sigma_{X^3}^2
    + 2\beta_3\beta_4 \Cov(X_i^2, X_i^3)
    > \sigma_u^2
  \label{eq:var_error_red}
\end{equation}

Al igual que en ejercicios anteriores, la varianza del error aumenta al
omitir variables relevantes, deteriorando la precision de toda la inferencia.

\begin{clave}
  \textbf{Respuesta al inciso~(c).} Si se eliminan $X_i^2$ y $X_i^3$:
  \begin{enumerate}[leftmargin=1.4em, itemsep=4pt]
    \item El coeficiente de $X_i$ resulta sesgado:
      $\E[\hat\alpha_2] = \beta_2 + \beta_3 b_{2,X^2} + \beta_4 b_{2,X^3}$.
    \item El sesgo combina efectos de signo opuesto ($\beta_3 b_{2,X^2} < 0$
      y $\beta_4 b_{2,X^3} > 0$), cuyo resultado neto es indeterminado sin
      datos concretos.
    \item $\hat\alpha_2$ pierde toda interpretacion economica como costo
      marginal en $X = 0$.
    \item El modelo lineal impone un costo marginal constante, contradiciendo
      la teoria.
    \item La varianza del error aumenta, invalidando la inferencia.
  \end{enumerate}
\end{clave}

\section{Sintesis comparativa}

\begin{table}[ht]
\centering
\caption{Consecuencias de mantener vs. eliminar los terminos polinomiales}
\label{tab:comparacion}
\small
\begin{tabular}{p{3.5cm} p{5cm} p{5cm}}
\toprule
\textbf{Criterio}
  & \textbf{Modelo cubico completo}
  & \textbf{Modelo lineal (omitiendo $X^2$, $X^3$)} \\
\midrule
Sesgo de coeficientes
  & Ninguno (insesgado bajo supuestos clasicos)
  & Sesgo $= \beta_3 b_{2,X^2} + \beta_4 b_{2,X^3}$ \\[4pt]
Interpretacion de $\hat\beta_2$
  & Pendiente del $CT$ en $X=0$ (costo marginal inicial)
  & Mezcla indistinguible de efectos lineal, cuadratico y cubico \\[4pt]
Forma del costo marginal
  & Cuadratica (forma de U): teoricamente correcta
  & Constante: teoricamente incorrecta \\[4pt]
Errores estandar
  & Inflados por colinealidad entre $X$, $X^2$, $X^3$
  & Menores, pero basados en un modelo incorrecto \\[4pt]
Validez de inferencia
  & Valida si los supuestos de MCO se cumplen
  & Invalida por omision de variables relevantes \\
\bottomrule
\end{tabular}
\end{table}

\appendix

\section{Ortogonalizacion de polinomios: la solucion correcta}

Si la colinealidad entre $X_i$, $X_i^2$, $X_i^3$ genera errores estandar
excesivamente grandes, la solucion correcta es construir polinomios ortogonales.
Para una muestra $\{X_1, \ldots, X_n\}$, los polinomios ortogonales se definen
recursivamente:

\begin{align}
  P_0(X_i) &= 1
  \label{eq:P0}\\
  P_1(X_i) &= X_i - \bar X
  \label{eq:P1}\\
  P_2(X_i) &= (X_i - \bar X)^2 - \frac{\sum (X_i - \bar X)^2}{n}
  \label{eq:P2}\\
  P_3(X_i) &= P_1(X_i)P_2(X_i) - \frac{\sum P_1(X_i)P_2(X_i)^2}{\sum P_1(X_i)^2} P_1(X_i)
  \label{eq:P3}
\end{align}

Por construccion, $\sum_i P_j(X_i) P_k(X_i) = 0$ para $j \neq k$. El modelo
estimado con estos polinomios:

\begin{equation}
  Y_i = \gamma_0 P_0(X_i) + \gamma_1 P_1(X_i) + \gamma_2 P_2(X_i) + \gamma_3 P_3(X_i) + u_i
  \label{eq:modelo_ortog}
\end{equation}

es algebraicamente equivalente al modelo cubico original pero tiene regresores
ortogonales, por lo que los errores estandar son minimos y cada $\hat\gamma_k$
se puede probar independientemente sin interferencia de los demas terminos.
Los coeficientes originales $\beta_2, \beta_3, \beta_4$ se recuperan mediante
la transformacion inversa.

\section{Condicion necesaria para la forma cubica del costo}

Desde el punto de vista de la teoria del costo, el modelo cubico~(\ref{eq:cubica})
implica que la funcion de costo marginal es:

\begin{equation}
  \text{CMg}(X) = \frac{dY}{dX} = \beta_2 + 2\beta_3 X + 3\beta_4 X^2
  \label{eq:CMg}
\end{equation}

Para que el costo marginal tenga forma de U (minimo en $X^* > 0$):

\begin{equation}
  \frac{d(\text{CMg})}{dX} = 2\beta_3 + 6\beta_4 X = 0
  \quad \Rightarrow \quad
  X^* = -\frac{\beta_3}{3\beta_4} > 0
  \label{eq:minimo_CMg}
\end{equation}

Esto requiere $\beta_3 < 0$ y $\beta_4 > 0$. Si se elimina $X_i^3$ (forzando
$\beta_4 = 0$), la condicion de minimo desaparece y el costo marginal no puede
tener forma de U, lo cual invalida la teoria.

\vspace{14pt}
\noindent\rule{\linewidth}{0.4pt}

\noindent{\small\textit{Nota.} Los valores de la matriz de correlaciones
provienen del ejemplo~7.4 de Gujarati~\& Porter (2010). La formula del sesgo
$\E[\hat\alpha_2] = \beta_2 + \beta_3 b_{2,X^2} + \beta_4 b_{2,X^3}$ es exacta
en muestras finitas bajo los supuestos clasicos de MCO.}

\end{document}
